\documentclass[11pt,a4paper]{article}
\usepackage[T1]{fontenc}
\usepackage{lmodern,amsmath,amssymb,amsthm,mathtools,booktabs,microtype}
\usepackage[margin=27mm]{geometry}
\usepackage[colorlinks=true,linkcolor=blue,citecolor=blue,urlcolor=blue]{hyperref}
\newtheorem{theorem}{Theorem}[section]
\newtheorem{lemma}[theorem]{Lemma}
\newtheorem{proposition}[theorem]{Proposition}
\newtheorem{corollary}[theorem]{Corollary}
\newcommand{\Q}{\mathbb Q}\newcommand{\Z}{\mathbb Z}
\DeclareMathOperator{\Tr}{Tr}\DeclareMathOperator{\rad}{rad}
\setlength{\emergencystretch}{2em}
\title{Thirteen Variables from a Weak $abc$ Inequality}
\author{}\date{9 November 2026}
\begin{document}\maketitle
\begin{abstract}
Assume the weak rational $abc$ inequality with exponent $121/100$,
together with the specified elliptic group, pole-realization and
simultaneous prime-pattern inputs. We construct a family of integer
polynomials in thirteen total variables. Every sufficiently large
fixed odd input shift gives rational input projection exactly $\Z$.
Each member has total degree less than $2^{44}$. A degree-four
rational map gives a logarithmic bound for an odd contact radical
from the weak inequality. An eighth-power height comparison then
bounds the rational input and one elliptic index together. Only
three elliptic points and two pole certificates occur, with six
local tests of order one. The coefficient threshold is ineffective.
Both the rational $abc$ conjecture and the ordinary Szpiro
conjecture imply the weak inequality used here.
\end{abstract}

\section{Statement and hypotheses}

All variable counts include the input. For a rational number $a/b$
in lowest terms with $b>0$, write
$h(a/b)=\log\max(|a|,b)$. For a nonzero integer, its radical
is the product of its distinct prime divisors, with positive sign.
Put $\kappa=121/100$. The additional height hypothesis used in this
paper is the following single inequality: there is a constant
$K_\kappa\ge1$ such that
\begin{equation}\label{eq:weak-abc}
 c\le K_\kappa\operatorname{rad}(abc)^{\kappa}
 \quad(a+b=c,\ a,b,c\text{ positive and pairwise coprime}).
\end{equation}
The same inequality, with the largest absolute value on the left,
applies to a coprime nonzero signed triple of sum zero, by moving
its negative terms to the other side.

Set
\[
 K=\Q(\sqrt2),\quad D=(3-\sqrt2)^2,\quad
 E:y^2=x^3-Dx,
\]
\[
 P=\left(\frac{57}{8}-\frac72\sqrt2,
                 -\frac{165}{8}+\frac{417}{32}\sqrt2\right).
\]
In addition to \eqref{eq:weak-abc}, retain the following arithmetic inputs with
exactly their stated scope in~\cite{Dependencies,Source27}.
\begin{enumerate}
\item $2E(K)=\Z P$, with the indicated non-torsion generator
and the source index-height bound. The direct point chart below
therefore describes the nonzero multiples of $P$.
\item The pole system, for $b\ne0$ and $B_0=17b$,
\[
 p^2=B_0(1-e)(1+3e),\qquad
 B_0q^2=(s^2-B_0e(B_0-b^2))(s^2+3B_0e(B_0-b^2)),
\]
excludes negative odd valuations of $b$ outside $\{2,3,17\}$.
For a prepared input
\begin{equation}\label{eq:prepared}
 \begin{gathered}
 0<b<17,\quad b\text{ has square reduced denominator},\quad
 b\in\Q_2^{\times2}\cap\Q_3^{\times2}\cap\Q_{17}^{\times2},\\
 v_2(b)\ge3,\quad v_3(b)\ge1,\quad v_{17}(b)\ge2,
 \end{gathered}
\end{equation}
the retained restricted-converse and descent hypotheses supply a
witness with $pqs\ne0$. The nonzero-witness consequence and its
single conic chart are proved in~\cite[Lemma~2]{Source27}.
\item The simultaneous prime-pattern input supplies the specified
prime contact numerators, Gaussian congruences, conjugate reduction
orders, fixed divisibility and avoidance, and compatible real
approximation for a fixed finite list of rational contact coefficients.
Its precise quantified version and the rational-contact and fixed-avoidance
consequences of Sun's theorem are given in~\cite{Dependencies}.
The even number of contacts below permits arbitrary positive rational
scales: Section~\ref{sec:complete} checks the square factors directly.
The order of choices is respected there.
\end{enumerate}
The exponent-nineteen contact-height input of the earlier source is
not needed here. It is replaced by Lemma~\ref{lem:contact}, which
uses \eqref{eq:weak-abc}.

The proved one-point case of the aggregate trace-height theorem
\cite{Trace} gives constants $A>0$, $J\ge0$, depending only on
$E,P$, such that
\begin{equation}\label{eq:traceheight}
 n^2\le A h(T(n))+J,\qquad
 T(n)=\Tr_{K/\Q}x([n]P)\quad(n\in\Z\setminus\{0\}).
\end{equation}
This applies von Buhren's theorem after checking the absence of
constant abelian translates for the fixed trace function. Its additive
constant is ineffective. It is a proved height theorem, not an
additional conjecture. For a nonzero multiple of $P$, the abscissa
is a positive square in both real embeddings; the curve equation
then gives $x>3\mp\sqrt2$ in the respective embeddings, hence
$T(n)>6$.

\begin{theorem}\label{thm:main}
Under the assumptions above, there is a computable family
\[
 F_B\in\Z[X,Y_1,\ldots,Y_{12}],\qquad B\ge3\text{ odd},
\]
with the following properties.
\begin{enumerate}
\item Every member is given by a complete finite integer coefficient
specification and satisfies
\[
             \deg F_B\le12339485119488<2^{44}.
\]
\item There exists a constant $B_*$ such that every odd $B\ge B_*$
has rational input projection exactly $\Z$.
\item No numerical sufficient value of $B_*$ is determined.
The supplied member $B=3$ is a coefficient-verification fixture,
not a polynomial certified to define $\Z$.
\end{enumerate}
The integer $B$ is a fixed coefficient parameter and is never a
quantified coordinate.
\end{theorem}

\begin{corollary}\label{cor:conjectures}
Retain the three source arithmetic inputs above. Either of the following
conjectures implies the conclusion of Theorem~\ref{thm:main}:
\begin{enumerate}
\item the ordinary rational $abc$ conjecture;
\item the ordinary Szpiro conjecture: for every $\epsilon>0$ there
is a constant $C_\epsilon$ such that
$|\Delta_E|\le C_\epsilon N_E^{6+\epsilon}$ for elliptic curves
$E/\Q$, where $\Delta_E$ is the minimal discriminant and $N_E$
the conductor.
\end{enumerate}
The semistable case of the second conjecture already suffices.
\end{corollary}
\begin{proof}
Rational $abc$ gives \eqref{eq:weak-abc} by taking
$1+\epsilon=121/100$. Oesterl\'e proves that the ordinary
Szpiro discriminant--conductor inequality implies the weak
$abc$ inequality with exponent $6/5+\epsilon$ for every
$\epsilon>0$ \cite[Section~I.3, p.~169]{Oesterle}.
Take $\epsilon=1/100$. The stated implication uses the
ordinary discriminant--conductor conjecture, not a separate
inequality involving both elliptic invariants $c_4$ and $c_6$.
\end{proof}


\section{Two consequences of the weak inequality}

Put
\[
 \Theta=\left\{0,-\frac12,-\frac14,-\frac25,-1,-\frac13\right\},
 \qquad S_0=\{p:p\le76176\}.
\]
For $t\in\Q\setminus\Theta$, let $M_0(t)$ be the product of
primes outside $S_0$ for which $v_p(t-\theta)$ is positive
and odd for some $\theta\in\Theta$.

\begin{lemma}\label{lem:contact}
Assume \eqref{eq:weak-abc}. There is a constant $C_0\ge0$ such that
\[
                    h(t)\le\frac{121}{74}\log M_0(t)+C_0
                  \qquad(t\in\Q\setminus\Theta).
\]
\end{lemma}
\begin{proof}
Consider the degree-four rational function
\[
 \varphi(z)=\frac{z(z+2)^3}{(2z+1)^3},\qquad
 \varphi(z)-1=\frac{(z-1)^3(z+1)}{(2z+1)^3}.
\]
Its complete inverse image of $\{0,1,\infty\}$ is
$\{0,-2,-1/2,1,-1,\infty\}$. For a primitive integer pair
$(x,y)$ away from these six points, put
\[
 A=x(x+2y)^3,\quad C=y(2x+y)^3,\quad
 A-C=(x-y)^3(x+y),
\]
\[
                G(x,y)=xy(x+2y)(2x+y)(x-y)(x+y).
\]
The polynomial identity for $A-C$ follows by expansion.
A prime dividing both $A$ and $C$ must divide both $x+2y$
and $2x+y$, so it is $3$. At that prime $x,y$ are units
and the smaller of those two valuations is one, because
$2(x+2y)-(2x+y)=3y$. Hence $\gcd(A,C)$ is either $1$ or $27$.
After division by this gcd, the three integers $A,C,A-C$
are pairwise coprime, and their radical is supported on $G$.

Write $H=\max(|x|,|y|)$. One has
$\max(|x+2y|,|2x+y|)\ge H$. If the corresponding prefactor
$|x|$ or $|y|$ is at least $H/4$, this gives
$\max(|A|,|C|)\ge H^4/4$. Otherwise the other prefactor
is $H$ and its linear factor is at least $H/2$, giving
$\max(|A|,|C|)\ge H^4/8$. Consequently
\[
                    h(\varphi(x/y))\ge4\log H-\log216.
\]
The weak inequality therefore implies
\[
 \log\operatorname{rad}(G(x,y))
       \ge\frac4\kappa\log H-\frac{\log216+\log K_\kappa}{\kappa}.
\]

For a nonzero integer $N$, let $\operatorname{sf}(N)$ be the
product of primes of odd valuation in $N$. Prime by prime,
$\operatorname{rad}(N)^2\le |N|\operatorname{sf}(N)$.
Also $|G(x,y)|\le36H^6$. Outside the fixed set $S_0$, two
of its six primitive linear factors cannot vanish together.
Thus its squarefree part outside $S_0$ is precisely the odd
contact radical for the six displayed points. The contribution
from $S_0$ is bounded by the fixed product of its primes.
Combining these estimates gives
\[
 \log M(z)\ge\left(\frac8\kappa-6\right)h(z)-C
             =\frac{74}{121}h(z)-C.
\]
Finally $z=2+1/t$ sends the six contacts in $\Theta$ to the
six points above. Its integral matrix has determinant $-1$;
its inverse is integral too, and both matrices have maximum
row sum three. Hence it changes logarithmic height by at most
$\log3$. Substitution of a primitive numerator-denominator pair
shows that the odd contact radicals agree outside $S_0$.
This proves the lemma after adjusting the fixed constant.
\end{proof}

For a rational pair $(u,v)$, define its projective height by
\[
 h_2(u,v)=h(1:u:v)=\log\max(D_0,|A'|,|B'|),
\]
where $D_0>0$, $u=A'/D_0$, $v=B'/D_0$ and
$\gcd(D_0,A',B')=1$. Both individual heights are at most $h_2$.

\begin{lemma}\label{lem:eighth}
Assume \eqref{eq:weak-abc}. There is a constant $C_1\ge0$ such that
\[
              h_2(u,v)\le\frac{121}{437}h(u^8+v^8)+C_1
                   \qquad(u,v\in\Q).
\]
One may take
$C_1=\frac{279}{1748}\log2+\frac{100}{437}\log K_\kappa$.
\end{lemma}
\begin{proof}
If one coordinate is zero, $h_2(u,v)=h(u^8+v^8)/8$, and
if both are zero all heights vanish. Otherwise use the primitive
representation above. Write $g=\gcd(|A'|,|B'|)$,
$A_0=A'/g$, $B_0=B'/g$. Then
$\gcd(A_0,B_0)=\gcd(g,D_0)=1$.
Put $H=\max(|A_0|,|B_0|)$ and $N=A_0^8+B_0^8$;
thus $H^8\le N\le2H^8$.

Let $k=\gcd(N,D_0^8)$, $m_0=N/k$, and $n=D_0^8/k$.
The fraction $u^8+v^8=m/n$ is reduced, where $m=g^8m_0$.
Put $M=\max(m,n)$, so the output height is $\log M$.
Every prime dividing $N$ either divides $D_0$ or survives
in $m_0$, giving $\operatorname{rad}(N)\le D_0m_0$.
Apply \eqref{eq:weak-abc} to the pairwise coprime positive
triple $A_0^8,B_0^8,N$:
\[
                H^{8/\kappa-2}\le K_\kappa^{1/\kappa}D_0m_0.
\]
Since $D_0^8=nN/m_0\le2nH^8/m_0$, it follows that
\[
 H^{437/121}\le2^{1/8}K_\kappa^{100/121}m_0^{7/8}n^{1/8}.
\]
Moreover $g=(m/m_0)^{1/8}$ gives
\[
               \max(gH,D_0)\le2^{1/8}H(M/m_0)^{1/8}.
\]
Taking logarithms and combining the last two inequalities yields
\begin{align*}
 h_2(u,v)
 &\le C_1+\frac{205}{1748}\log m_0
                 +\frac{121}{3496}\log n+\frac18\log M\\
 &\le C_1+\frac{121}{437}\log M.
\end{align*}
Here $m_0,n\le M$, all coefficients are nonnegative, and
$205/1748+121/3496+1/8=121/437$. This accounts for the
common numerator factor and every denominator cancellation.
\end{proof}

\section{The complete rational source}

Fix an odd integer $B\ge3$ and set
\begin{equation}\label{eq:input}
                            \beta=2X^2+B.
\end{equation}
For rational $X=u/v$ in lowest terms, integrality of $\beta$
implies $v^2\mid2$, hence $v=1$. Thus $\beta\in\Z$ if and
only if $X\in\Z$. For integer input $\beta$ is odd, and
$\beta\ge B$ for every rational input.

Retain three pairs of rational coordinates, one for each of the
points $A,H,U$. For any such pair $(h_0,h_1)$ put
\[
 h=h_0+h_1\sqrt2,\quad x=h^2,\quad j^2=h^4-D,\quad y=hj.
\]
The direct point-chart theorem of~\cite{Source27} proves that
these are exactly the nonzero multiples of $P$. Each finite
choice $j$ is encoded by one scalar monic quartic as follows.
Write $h^4-D=a+b\sqrt2$, put
$c=2a+b$, $d=-a-b$, and write
$(c+d\sqrt2)^2(a+b\sqrt2)=a'+b'\sqrt2$.
Then
\[
 Y^4-4a'Y^2+8(b')^2=0,\quad
 j=\frac{Y/2+(b'/Y)\sqrt2}{c+d\sqrt2}.
\]
Its discriminant is
$2^{15}(b')^2(2a^2-b^2)^4(a^2-2b^2)^2$, nonzero at every
rational pair $(h_0,h_1)$. Indeed $h^4-D\ne0$ by its norm,
the two quadratic norm factors are anisotropic over $\Q$, and
$b'=-4a^3+6ab^2+3b^3$ is anisotropic by Eisenstein at 3.
The square-root chart is complete because its transformed rational
component is $2(u^3-2v^3)$ for $j=u+v\sqrt2$, which is nonzero
for $j\ne0$. These are algebraic facts independent of the height hypothesis.

Compute $V_0=H-P$ and $V_1=U-A$ by the usual affine addition
formulas. Put $f=1/x$, $z=-x/y$, let $\sigma$ be conjugation
of $K/\Q$, and define
\[
 L=\frac{z(U)}{z(H)},\qquad
 L'=\sigma\!\left(\frac{z(V_1)}{z(V_0)}\right).
\]
Use precisely the six field-valued tests
\begin{equation}\label{eq:tests}
 D_1=f(H),\ D_2=f(U),\ D_3=\sigma f(V_0),\
 D_4=\sigma f(V_1),\ D_5=L-L',\ D_6=L-\beta.
\end{equation}
Every test is assigned order one. No scalar-height pole or Lucas
relation is present.

Let $T=\Tr x(A)$, and use the two positive rational scales
\[
                         a_0=1,\qquad a_1=\beta^8+T^8.
\]
Retain one rational slope $\lambda$. For $r=0,1$, set
\begin{align}
 t_r&=a_r\lambda,&\delta_r&=t_r-1,&
 F_r&=\prod_{\theta\in\Theta}(t_r-\theta),&
 R_r&=\frac{F_r}{\delta_r^8},\label{eq:detector}\\
 W_r&=\prod_{j=1}^6\left(1+
             \frac{D_j^2}{\delta_r^4R_r}\right),&
 b_r&=\frac{R_r}{\Tr(W_r^{-1})^2}.\label{eq:W}
\end{align}
All displayed divisions, affine-chart denominators, $F_r$,
$\delta_r$, $\operatorname{Norm}(W_r)$, and $\Tr(W_r^{-1})$
are guarded nonzero. These conditions are part of the source.

For each $b=b_r$, retain independent rational parameters $\tau_r,s_r$.
Put $B_0=17b$ and
\[
 e=\frac{\tau_r^2-B_0}{\tau_r^2+3B_0},\qquad
 q^2=\frac{(s_r^2-B_0e(B_0-b^2))
                  (s_r^2+3B_0e(B_0-b^2))}{B_0}.
\]
Guard the denominator and require $q\ne0$. This reconstructs
the first pole ordinate as
$p=-4B_0\tau_r/(\tau_r^2+3B_0)$.
Each rational-square equation $q^2=N_0/D_0$ is made monic
by $Z=D_0q$, giving $Z^2=N_0D_0$ with $Z\ne0$ and $D_0\ne0$.
No new unrestricted rational parameter is used for these roots.

\section{The local implication}

\begin{lemma}\label{lem:local}
For any rational solution of the guarded source, write
$A=[n]P$, $H=[h]P$, $U=[u]P$.
If $\beta\ne n$ and $N$ is the positive absolute numerator
of $\beta-n$ in lowest terms, then
\[
                         M_0(t_r)\mid N\qquad(r=0,1).
\]
\end{lemma}
\begin{proof}
Fix a prime $p\mid M_0(t_r)$ and abbreviate $R=R_r$, $W=W_r$.
Outside $S_0$ only one contact factor has positive valuation,
$\delta_r$ is a unit, and $v_p(R)$ is odd because its denominator is an eighth power. For each place $\mathfrak p$ above $p$, the two terms
of a factor of $W$ have valuations of different parity. Hence
\[
 v_{\mathfrak p}(W)=\sum_j
       \min\{0,2v_{\mathfrak p}(D_j)-v_p(R)\}\le0.
\]
Put $w=\max_{\mathfrak p\mid p}v_{\mathfrak p}(W)\le0$.
Then $v_p\Tr(W^{-1})\ge-w$.
The valuation of $b=R/\Tr(W^{-1})^2$ is odd, so the pole
obstruction gives
\[
                   0<v_p(b)\le v_p(R)+2w.
\]
Thus $\nu=v_p(R)>0$ and $w>-\nu/2$. At a place attaining
$w$, each summand in the valuation of $W$ is at least $w$.
It follows that $v_{\mathfrak p}(D_j)>\nu/4>0$, so all
six tests have valuation at least one at that place.

The first four tests put $H,U$ into the formal group at
$\mathfrak p$ and $V_0,V_1$ into it at the conjugate place.
At these good unramified places the formal logarithm satisfies
\[
                         \log_E(Q)=z(Q)(1+O(p^4)).
\]
Indeed its first omitted term is a fixed multiple of $Df(Q)^2$;
$v_p(f(Q))\ge2$ in the formal group. The remaining coefficient
estimates, uniformly for $p>76176$, are those proved in the
formal-logarithm paper~\cite{Logs}. The exact logarithmic ratios
are $u/h$ and $(u-n)/(h-1)$, respectively.

The affine guards imply $h\ne0,1$. At least one of $h,h-1$
is a $p$-adic unit, so the corresponding exact ratio is integral.
Relative precision four makes its truncated ratio integral.
The test $L-L'\equiv0\pmod p$ makes the other truncated ratio
integral as well; the same relative precision then makes the
other exact ratio integral. Relative errors are now absolute
errors of valuation at least four. Consequently
\[
 L\equiv u/h,\qquad L'\equiv(u-n)/(h-1),\qquad
 L\equiv L'\pmod p.
\]
Multiplying by the integral numbers $h,h-1$ and subtracting yields
$L\equiv n\pmod p$. The target test gives $L\equiv\beta\pmod p$.
In particular $\beta$ is $p$-integral, and $p$ divides the
reduced numerator of $\beta-n$. This holds at every consumed
prime and proves the divisibility. The same argument includes
inert primes; if their formal-group requirements are incompatible,
no source solution occurs at that prime.
\end{proof}

\section{Soundness and the uniform input shift}

Suppose a rational source solution has $\beta\ne n$. By
Lemmas~\ref{lem:local} and~\ref{lem:contact},
\[
 h(\lambda),h(a_1\lambda)\le\frac{121}{74}\log N+C_0.
\]
The quotient-height inequality gives
\[
 h(a_1)\le\frac{121}{37}\log N+2C_0.
\]
Apply Lemma~\ref{lem:eighth} to $\beta,T(n)$, and put
\[
 \eta=h_2(\beta,T(n)),\qquad
 q=\frac{14641}{16169}<1,\qquad
 C=\frac{242}{437}C_0+C_1.
\]
Then
\begin{equation}\label{eq:combined}
                         \eta\le q\log N+C.
\end{equation}
All these constants depend only on the fixed contact map,
the fixed power eight, and the weak-$abc$ constant. They are
independent of the shift $B$, the input, and every witness.

Writing $\beta=a/b$ in lowest terms gives
$N=|a-nb|\le e^{h(\beta)}(1+|n|)$, and $h(\beta)\le\eta$.
Consequently
\[
 \eta\le c\log(1+|n|)+C',\qquad
 c=\frac{q}{1-q},\quad C'=\frac{C}{1-q}.
\]
Since $h(T(n))\le\eta$, the trace-height inequality implies
\[
 |n|^2\le Ac\log(1+|n|)+AC'+J
           \le Ac|n|+AC'+J.
\]
Choose an integer $M>Ac+\sqrt{AC'+J}$. Then $|n|<M$, so
\[
                    h(\beta)\le\eta
                       \le c\log(1+M)+C'=:H_*.
\]
Choose $B_*>e^{H_*}$. For every fixed odd $B\ge B_*$,
\eqref{eq:input} gives $h(\beta)\ge\log\beta\ge\log B>H_*$,
a contradiction. Thus $\beta=n\in\Z$, and the elementary
input-map argument yields $X\in\Z$.
This proves soundness uniformly for every sufficiently large shift.
It does not determine that shift effectively.

\section{Completeness}\label{sec:complete}

Fix an integer input $X$ and an odd coefficient $B\ge3$.
Then $n=\beta$ is an odd integer at least three. Fix
$A=[n]P$ and its two positive scales before choosing
any prime. They are distinct because $a_1>1$.
Temporarily put $H=[h]P$ and $U=[hn]P$.

Choose a nonexceptional $h_*\in64\Z$, and work in the real
orbit closure of $64\Z P$. Choose a compact real neighborhood
for $H$ around $[h_*]P$ on which the point coordinates, both difference charts,
and the six tests are defined and bounded. Such a neighborhood
exists: the rational functions restricted to the multiplication
curve $H\mapsto(H,[n]H)$ are not identically singular.
The difference charts exclude $h=\pm1$, while every nonzero
multiple of the non-torsion $P$ has nonzero abscissa and ordinate.
Pass to a smaller compact neighborhood to avoid all remaining
zero denominators. Enlarge $L$ with $64\mid L$ and the real
component-group condition so that the progression
$h\equiv h_*\pmod L$
preserves fixed lower valuation bounds for these tests at
$2,3,17$. These are precisely the local-continuity and real-orbit
steps of the source realization, with fewer point roles and no
scalar constraint.

Choose a fixed $D_0$ with $D_0^2\theta/a_r\in\Z$ for all
contacts and both scales, set $v=D_0^2w^2$, and seek
$\lambda=u/v$. Apply the retained prime-pattern theorem to
the distinct forms $u-(\theta/a_r)v$, after repetitions have
been removed. All bad primes and previously fixed moduli are
excluded before this application. The theorem permits arbitrarily
large divisibility of $w$, a sufficiently large prescribed lower
bound on $u/w^2$, and the conjugate reduction-order conditions.

As in~\cite{Dependencies,Baseline}, the annihilating order moduli
$U_0,V_0$ for the two conjugate families satisfy
\[
 \gcd(U_0,V_0)=1,\qquad \gcd(U_0,L)=16,\qquad\gcd(V_0,L)=1.
\]
Thus the congruences
$h\equiv0\pmod{U_0}$, $h\equiv1\pmod{V_0}$,
$h\equiv h_*\pmod L$ are compatible, since $16\mid h_*$.
Choose $h$ in that
progression while retaining the real neighborhood and avoiding
the finitely many excluded multiples. At one place above each
contact prime, $H,U$ are formal; at the conjugate place,
$H-P,U-A$ are formal. Both exact logarithmic ratios are $n$.
The constant approximation has precision four, so all six tests
have their assigned order one, including the target $n=\beta$.

The inherited detector preparation now applies without alteration.
For clarity, its main ingredients are the following. The identity
\[
 R_r=\frac{a_r^6\prod_{\theta\in\Theta}(u-(\theta/a_r)v)}
                  {v^6(a_r\lambda-1)^8}
\]
has a rational-square prefactor. In particular, the possibly nonsquare
scale contributes the square $a_r^6$.
At each chosen contact prime $v_p(R_r)=1$; the successful
conjugate place gives all inverse norm-factor residues equal to
one, and the trace of $W_r^{-1}$ is a unit with residue one or two.
The resulting $b_r$ has square reduced denominator, and the
prescribed congruences give its local squareclasses at $2,3,17$.
The fixed lower bounds on the tests and sufficient divisibility
of $w$ give the depths in~\eqref{eq:prepared}.
On the compact real set, uniformly,
\[
 R_r\sim a_r^{-2}\lambda^{-2},\qquad
 W_r=1+O(\lambda^{-2}),\qquad
 b_r\sim(4a_r^2\lambda^2)^{-1}.
\]
Thus $0<b_r<17$ for the sufficiently large lower bound on
$\lambda$ fixed before prime selection. All detector guards
are nonzero. These calculations are the source comparison-pole
preparation proof~\cite{Baseline}; they depend on the finite
list of tests and their orders, not on a Lucas or scalar test.

Each prepared input has an independent pole certificate with
$pqs\ne0$ by the retained pole-realization hypotheses.
Since $p\ne0$ excludes $e=1$, the substitution
$\tau=p/(e-1)$ gives the single chart used above and its
nonzero denominator. Thus both nonzero quadratic-root conditions
are realized. The direct point charts cover $A,H,U$.
Every rational input $X\in\Z$ therefore has all required
witnesses, for every fixed odd $B\ge3$.

\section{The polynomial, its bounds, and verification}

The retained coordinates are exactly
\begin{center}
\begin{tabular}{lr}
\toprule
Coordinates & Number\\
\midrule
Input $X$ & 1\\
Three point pairs for $A,H,U$ & 6\\
Comparison slope $\lambda$ & 1\\
Two independent pole pairs $(\tau_r,s_r)$ & 4\\
\midrule
Retained before finite-root elimination & 12\\
\bottomrule
\end{tabular}
\end{center}
There are three monic point quartics and two monic pole quadratics.
The point algebra has rank $4^3=64$, and the full algebra has
rank $4^3 2^2=256$. Its complete guard is $gq_1q_2$ with
$g$ in the point algebra, and each quadratic radicand also belongs
to that algebra. There are no further residual equations.
The separating-element theorem~\cite{Active} therefore replaces
all five roots and all nonzero guards by one rational coordinate,
without changing the input projection. This gives thirteen total
coordinates. Its determinant formulas and its 64-factor product,
applied to the supplied integer coefficient circuit, define $F_B$
completely.

The native generator weights are
\[
                         6,6,6,51131,69765.
\]
The full guard and its point factor have weights $148317$ and
$27421$. In the notation of~\cite[Theorem~4.1]{Active},
the multiplication-element weight is at most $W_1r+W_0$, where
\[
                    W_1=9577692,\qquad W_0=70636907.
\]
Thus
\[
 \deg F_B\le256\bigl(9577692\cdot64^2
                         +(70636907+1)\cdot127\bigr)
             =12339485119488<2^{44}.
\]
The input shift is a named coefficient parameter of weight zero
in the coefficient graph. No operation or branch depends on its
numerical value. Consequently these weight inequalities, and hence
the degree bound, hold symbolically for every fixed $B$.

The directory \nolinkurl{computation/few_variables/weak_abc13}
contains the rational-source builder, its integer coefficient
specification, and independent source and bound certificates.
The source has 1758 rational-circuit nodes and its integer
coefficient graph has 754 nodes. The independent conversion auditor
checks every affine substitution, the three point quartics and
their automatic nonzero factors, both pole-square clearings,
and all 48 atomic guard factors by exact expression identities.
The separate finite-algebra checker verifies coefficient dependencies,
companion-column inequalities and all 64 degree cases without
importing the builder.

For the algebraic fixture $B=3$ it certifies
\[
 \log_2 H(F_3)\le29454018907833<2^{45}.
\]
This is not a height bound for a member certified to define $\Z$.
The proof provides an ineffective sufficient shift, whereas the
coefficient algorithms construct and bound each specified member.
The exact finite checks establish their recorded algebraic assertions;
they do not verify \eqref{eq:weak-abc} or the inherited arithmetic premises.

\begin{thebibliography}{9}
\bibitem{Oesterle}
J.~Oesterl\'e, Nouvelles approches du ``th\'eor\`eme'' de Fermat,
\emph{Ast\'erisque} \textbf{161--162} (1988), S\'eminaire
Bourbaki, expos\'e no.~694, 165--186; Section~I.3, p.~169.
\url{https://www.numdam.org/article/SB_1987-1988__30__165_0.pdf}.
\bibitem{Source27}
\emph{A defining polynomial in twenty-seven variables:
an existence theorem with a uniform degree bound},
companion paper, \nolinkurl{papers/few_variables27}, 9 November 2026.
\bibitem{Dependencies}
\emph{Arithmetic dependencies of the construction with 29 variables},
companion paper, \nolinkurl{papers/arithmetic_dependencies}.
It pins the source statements and Sun's prime-pattern theorem.
\bibitem{Baseline}
\emph{An explicit polynomial in 29 variables over the rationals},
companion paper, \nolinkurl{few_variables_20261108_v8.tex},
local implications and simultaneous-completeness sections.
\bibitem{Trace}
\emph{Aggregate trace heights of elliptic multiples}, companion
paper, \nolinkurl{papers/trace_heights}, applying J.~von Buhren,
\emph{J. Th\'eor. Nombres Bordeaux} \textbf{29} (2017), 307--320,
Theorem~1.1, \url{https://www.numdam.org/item/10.5802/jtnb.981.pdf}.
\bibitem{Logs}
\emph{Uniform formal-logarithm estimates}, companion paper,
\nolinkurl{papers/elliptic_logarithms}.
\bibitem{Active}
\emph{Separating the Nonzero Part of a Finite Algebra},
companion paper, \nolinkurl{active_finite_algebra_20261109_v15.tex}.
\end{thebibliography}
\end{document}
